\section{Discussion: Cross-Domain Transfer and the DeFi-as-HFT Hinge}
\label{sec:discussion}

\Cref{sec:empirical} reported the binding result of the pre-registered
three-hypothesis matrix.  This section steps back to position the
methodology within the high-frequency-trading (HFT) microstructure
canon, identify the structural reason the event-time formulation is
the right resolution for DeFi lending, and bound the claim with the
limitations that the four-month $T$ and the absence of MEV protection
in the backtest enforce.

The central conceptual move of this paper is the recognition that the
on-chain analog of MacKenzie's regression-based ``squashing'' of
microstructure signals \cite[p.\,176]{mackenzie2021} ---
$X^{T} X$ as the literal XTX Markets nameplate --- is realised
exactly by a multi-criteria decision allocator over rate,
utilization, gas-cost, and stability features.  Both reduce a
high-dimensional, heterogeneous-quality signal vector to a
one-dimensional act: choose-venue in HFT, choose-protocol in DeFi.
This is the methodology bridge that the architectural-transfer
hypothesis of our earlier LOB-to-DeFi work
\cite{solovev2026lob,solovev2026vault} could not directly verify;
event-time evaluation closes it, and the multicriteria precedent in
crypto-portfolio selection \cite{aljinovic2021} supplies the
weighting-scheme idiom.

\subsection{Signal-class taxonomy and the F1--F4 mapping}
\label{sec:disc-signals}

\cite[Table 3.2, p.\,97]{mackenzie2021} taxonomises HFT signals into
four classes: \emph{futures lead}, \emph{order-book dynamics},
\emph{fragmentation}, and \emph{related instruments}.  We claim that
each maps to a concrete, in-principle measurable DeFi-lending feature
family:
\begin{description}
  \item[F1 (lead).] The Maker DSR redemption rate, the sDAI proxy,
    and the Curve 3pool USDC/USDT/DAI swap-rate move ahead of Aave
    and Compound USDC supply rates with a lag of roughly six hours,
    by the elasticity-of-substitution mechanism that
    \cite{krause2005} documents for closed-form market-depth
    substitution in unified-pool systems.  Source: subgraph events,
    free.
  \item[F2 (order-book dynamics).] Pending mempool transactions
    targeting each pool --- deposits, withdrawals, borrows, repays
    --- are observable \emph{before} they execute, exactly as a
    Glosten--Milgrom dealer observes a quote-revision intent before
    the trade \cite[Ch.\,3]{ohara1995}.  This is the most direct
    analog of the ATD ``Goldman leaves its bid'' signal that
    \cite{mackenzie2021} describes on pp.\,102--104.  We defer F2 to
    future work because historic Flashbots mempool snapshots have
    documented gaps in 2024--2025 (design-spec risk R1).
  \item[F3 (fragmentation).] The cross-protocol rate spread itself
    --- $15$ pairwise spreads for the six-way panel --- and the
    cross-protocol utilization spread.  This \textbf{is} our decision
    variable, and \Cref{sec:emp-ablation} confirms it as the
    dominant contributor to T3's hazard fit, with the synthetic-Cox
    \texttt{f3\_spread\_max\_minus\_min} coefficient at $\approx
    +500$ vs noise-level F1/F4 coefficients.
    \cite{gudgeon2020loanable} reports a Compound-leads-Aave
    cointegration with a $0.607$ speed-of-adjustment coefficient;
    our event-time replay shows that this leads-and-lags structure
    is exploitable at the per-block granularity but not at the
    hourly aggregation level used in our prior published work
    \cite{solovev2026vault}, which is the empirical resolution of
    that earlier $H_{0}$ verdict.
  \item[F4 (related instruments).] ETH spot price (drives gas
    regime), top-LP concentration per pool (more transparent
    on-chain than in MacKenzie's anonymous order books), and USDC /
    USDT / DAI peg deviations as leading indicators of liquidity
    stress; the Aave-vs-Compound cross-chain comparative work of
    \cite{aavecompoundcrosschain2025} catalogues the relevant
    risk-parameter levers.
\end{description}
The MacKenzie-to-DeFi mapping is the paper's central
\emph{theoretical} contribution and stands independently of the
$H_{1}$ binding outcomes.  The Kissell impact-vs-edge inequality
\cite{kissell2014} supplies the gating threshold that converts each
of F1, F3, and F4 into an actionable trigger.

\subsection{The hinge: allocator + protocol launches as mutually reinforcing}
\label{sec:disc-hinge}

Andrew Abbott's notion of a \emph{hinge}, recapitulated by
\cite[pp.\,93--94]{mackenzie2021}, describes a process that creates
rewards in more than one sphere of activity --- in MacKenzie's case
the Island ECN's profits funded the HFT firms that traded on Island,
and those firms' liquidity, in turn, gave Island its share-of-volume
victory over rival venues.  The mutual reinforcement closes a
feedback loop that neither side could close alone.

The DeFi analog is sharper than the HFT original.  A multi-protocol
gas-aware allocator's fragmentation-rent (arbitrage across Aave,
Morpho, Spark, Compound, Fluid, Euler rate quotes) simultaneously
\textbf{(a)} earns the allocator alpha, and \textbf{(b)} makes it
rational for a new lending protocol to launch --- because
differentiated-rate venues now attract sophisticated flow that
bootstraps TVL, which feeds back into the rate-quote signal that the
allocator profits from \cite{gudgeon2020loanable}.  The empirical
evidence is the protocol ordering visible in \Cref{tab:test-matrix}:
Spark, Morpho Blue, Fluid, and Euler V2 all launched after the
2024-Q4 baseline window, and each is meaningfully on the Pareto
frontier of (yield, TVL, switching cost) for at least one of the
test-window quarters reported in \Cref{tab:regime-breakdown}.  The
pre-condition that \cite[p.\,95]{mackenzie2021} requires for a hinge
to close --- \emph{unified clearing} --- DeFi enjoys for free,
because Ethereum L1 is the shared settlement substrate for all six
in-scope protocols.  Cross-chain allocation between Solana and
Ethereum would not close the hinge in this sense, because the
asynchronous bridges break the unified-clearing pre-condition.

This recasts the paper's contribution from ``a faster MCDM allocator''
to ``a measurement of the hinge''.  The Sharpe gap that T1 captures
over B4 --- $+61$\,bp net APY on the Sep--Dec 2025 validation slice
(commit \texttt{398809d}), persisting into the test window per
\Cref{sec:emp-results} --- is not just engineering; it is the rent
that any allocator at the per-block resolution can extract from the
multi-venue equilibrium, and its sign and magnitude are an
observable about the maturity of the DeFi-lending market.  The
gas-cost-crossover finding of Plan~C (at the calibration point of
\$$17.5$ gas per rebalance on a \$1\,M position, a switch pays only
if the spread is $\ge 5$\,pp or the position is $\ge 5\times$ scale)
quantifies the minimum hinge-rent visible to any participant: below
that floor, the allocator does not close the loop and protocol
launches do not bootstrap.  Above that floor --- exactly the regime
the gas-aware T1 selects into --- the hinge is open.

\subsection{Limitations: MEV exposure and the asymmetric-speed-bump fix}
\label{sec:disc-limits}

The backtest is honest about one limitation: it does not deduct
miner-extractable-value (MEV) losses on the rebalance transaction.
With public-mempool submission, a \$$1$\,M rebalance would face an
expected $5$--$30$\,bp sandwich-extraction tax, which would erode
much of the T1 edge documented in \Cref{tab:test-matrix}.  The
published agent therefore submits rebalances through the Flashbots
private mempool via \texttt{eth\_sendPrivateTransaction}, as
described in \Cref{sec:methodology}.

We reframe this protection in MacKenzie's terms.
\cite[pp.\,200--203]{mackenzie2021} characterises IEX's symmetric
$350$-microsecond coil --- the speed bump that delays \emph{all}
order types equally --- as a partial solution: it protects slow
takers at the cost of also slowing market-maker quote cancellations.
The Flashbots private mempool is structurally different: it delays
\emph{visibility} to MEV bots until inclusion --- an asymmetric speed bump
in MacKenzie's sense --- while our own intent-cancellation and re-pricing remain
fast.  The closer FX analog is ``last look'' rather than the IEX
coil; we adopt that framing in the methodology section.  The Plan E
live-agent task verifies the Flashbots submission path end-to-end on
Sepolia testnet, completing the bridge from backtest to production
that the $H_{1}$ binding tests presuppose, and the
\cite{aavecompoundcrosschain2025} cross-chain comparison establishes
the parameter-symmetry conditions under which this generalises.

The second limitation, addressable in future work, is the $T = 4$
monthly Sharpe sample size.  With $T = 12$ (one calendar year), the
$\sqrt{T-1}$ pre-factor in the Lopez de Prado DSR formula
\cite[Ch.\,14.7.3]{lopezdeprado2018} would relax by a factor of
$\sqrt{11/3} \approx 1.92$ and the gate would become substantially
easier to clear; the conservative posture of the Plan~D test window
is deliberate and traceable to the panel's coverage start at
2024-11-01 and the strict purge-and-embargo requirement of AFML
Ch.~7.4.  The third limitation is the deferred F2 mempool channel:
once Flashbots historic snapshots stabilise, the
$X^{T} X$-style regression of \cite[p.\,176]{mackenzie2021} can be
augmented with the order-book-dynamics column of MacKenzie's
Table~3.2, which we expect will close a substantial fraction of the
gap to the Kissell information ceiling \cite{kissell2014} ---
specifically the $\sim 40\%$ residual that \Cref{sec:emp-ablation}
attributes to unforecastable adverse-selection noise per
\cite{ohara1995}.

The honest-$H_{0}$ posture closes the discussion.  The publication
protocol commits to reporting the binding result whether or not
$H_{1}$ is rejected, and the methodology contribution --- event-time
decision frame, F1/F3/F4 signal taxonomy, hinge measurement, and
production-grade fetcher infrastructure --- stands either way.  The
allocator-as-XTX bridge to \cite{mackenzie2021} and the
hinge-measurement reframing of the contribution are the durable
deliverables; the H1 verdict is an empirical update on the maturity
of the multi-protocol equilibrium at the time of the test window.

% end of section
